% Pursuing Stacks (cotes 134-2 to 134-8): archive cote 134-7, site page 61 (PDF page 62), typescript page 503.
% typescript. Draft compared: ps6.tex:2992-3051.
% AI-transcribed by Claude Opus 5.5 (claude-opus-5-5, adaptive thinking, effort high), Message Batches API, batch ps4 (msgbatch_012Aoc1NuskATPsyBeKSYK7J), 2026-11-12; not yet checked by eye.
\scanpage{134-7/61}
\mspage{503}

Maybe it is worthwhile still to give a down-to-earth formulation of
such a comparison statement. For any free $\underline{Z}$-module $M$\inserted{of finite type}, let
\[ \mathrm{Symbin}^{\hand{\unsure{*}}}_{\underline{Z}}(M) \hand{\subset} \mathrm{Sym}^{*}_{\underline{Q}}(M_{\underline{Q}}) \qquad \text{(where } M_{\underline{Q}} = M\otimes_{\underline{Z}}\underline{Q} \text{ )} \]
be the subalgebra of the algebra of polynomial functions on $M_{\underline{Q}}^{\hand{\vee}}\hand{\simeq}(M^{\hand{\vee}})_{\underline{Q}}$
which are \unsure{integgal|integral} - valued on $M^{\hand{\vee}}$ = dual module of M. Now let $L_{*}$
be any semisimplicial $\underline{Z}$-module whose components are free of finite type,
and consider the canonical map of cosemisimplicial $\underline{Z}$-modules
\[ \mathrm{Symbin}_{\underline{Z}}(L_{*}^{\hand{\vee}}) \hand{\to} \mathrm{Maps}(L_{*}, \underline{Z}) \quad , \]
described componentwise in an obvious way. The question is whether this
is a weak equivalence i.e. induces a \xout{chain equiva} quasi-isomorphism for
the associated cochain complexes, under the extra assumption that $L_{*}$ is
0-connected i.e. the associated chain complex has zero $H_{0}$ (and possibly,
if needed, assuming even 1-connectedness i.e. $H_{0}$ \uline{and} $H_{1}$ of the associated
chain complex are \inserted{both} zero). Presumably, by easy devissage arguments one
should be able to reduce to the case when $L_{*}$ is a $K(\underline{Z},n)$ type, \xout{and} and
more specifically still, that it is the semisimplicial abelian complex
associated to the chain complex reduced to $\underline{Z}$ placed in degree n (where
$n\hand{\geq}1$). Thus, the question is whether Eilenberg-Maclane cohomology (with
coefficients in $\underline{Z}$) for $K(\hand{\pi},n)$ types (or more specifically, $K(\underline{Z},n)$ types)
can be \xout{computed}\inserted[typed]{expressed} in terms of derived functors of the $\mathrm{Symbin}(M^{\hand{\vee}})$ functor .
At any rate, whether Symbin is just the right functor to fit in or not,
it looks like an interesting question whether Eilenberg-Maclane cohomology
(or, more relevant still, homology) can be expressed in ter\over{\unsure{n}}{m}s of the deri-%
ved functor of a suitable non-additive contravariant (resp. covariant)
functor\inserted{B} from (Ab) (or from the subcategory of free $\underline{Z}$-modules) to itself.
If so, there should be a \xout{(possibly somewhat sophisticated)} way of defi-%
ning a (possibly somewhat sophisticated) notion of ``B-schematic homotopy
types'' (over a ground ring k endowed with suitable extra structure, such
as a binomial structure ), \inserted{in terms of ``unipotent B-bundles'',} in such a way that any \xout{\illegible{1}} ``discrete'' namely
usual homotopy type, satisfying a suitable 1-connectedness restriction,
admits an essentially unique ``B-schematization''.

I do'nt feel like pursuing these questions here, which would take me
too far off the main line of investigation I've been out for. At any rate,
whether or not Eilenberg-Maclane \xout{co}homology may be expressed in terms of
the \inserted{total} \xout{xxxxx}\inserted[typed]{left} derived functor of a suitable functor from (Ab) to itself, it
would seem that the somewhat naive approach towards schematic homotopy
types we have been following, valid over an arbitrary (commutative)
ground ring k without any extra structure needed on it, is worthwhile

\margin{61}
